Constraint-based models such as flux balance analysis (FBA) predict cellular
metabolic states by solving linear optimization problems. When
nutrients or enzyme capacities vary continuously, optimal fluxes do
not change smoothly but follow piecewise-linear trajectories, changing slope
abruptly at bottlenecks. Here we show that the response's true
``curvature'' is not an ordinary function but a
discrete measure concentrated on the boundary interfaces where active
constraints switch. Path-dependent memory (holonomy) scales linearly with
perturbation size (slope $1.00$; smooth systems quadratic), and
$93.4$--$100.0\%$ of second-order adjustment concentrates at discrete
bottleneck transitions. These switches bridge to smooth
geometry: under refinement they converge weakly to
curvature densities and strongly in $L^1$, but total variation
fails to converge generically, defining the resolution window ($h \ll
\sigma \ll L_{\mathrm{var}}$).

Integrating this measure along physiological paths yields a
parameter-free gene sensitivity metric, $\kmu$: each enzyme's
rerouting burden. In carbon-starved \emph{E.~coli}, $\kmu$
predicts transcriptional induction in $424$ genes
($r = +0.395$, $p = 2.6 \times 10^{-17}$),
robust across tie-breaking rules and stress axes. Induced
genes sit in operons of the global carbon and energy regulons;
the rerouting mass concentrates on the fork metabolites of
central carbon. The association vanishes at the protein layer
($r = -0.083$, $366$ genes, matched proteomics):
cells transcribe standby capacity for rerouting, buffering protein
synthesis, so cyclic-perturbation memory ($66\%$ non-reverting)
must live in fast post-translational state.
Double-knockout epistasis mirrors active-set boundary overlaps
($\rho_S = 0.865$). Our framework unites linear programming,
discrete geometry, and transcriptional regulation into an
accessible, predictive foundation for metabolic systems biology.
